\documentclass{letter}

% \usepackage[superscript,biblabel]{cite}
% \bibliographystyle{naturemag}

\signature{Jerome Kelleher}

\address{Big Data Institute\\University of Oxford\\UK}
\begin{document}

\begin{letter}{Nature Methods}

\opening{Dear Dr Korda,}

We thank you and the reviewers for the careful and constructive evaluation of
our Correspondence, ``Population-scale Ancestral Recombination Graphs with
tskit 1.0.'' We are grateful for the positive overall assessment and for the
specific suggestions, which have helped us improve the manuscript.

In response, we have revised both the main text and the Supplementary
Information. In particular, we have clarified the significance of the tskit 1.0
release, clarified how the present paper differs from our earlier publications,
added a new supplementary section (``Tskit functionality and benchmarks'')
together with Table S2, and added a brief statement in the main text on likely
future development directions.

We have included a copy of the manuscript in which all changes have been
highlighted.

Below we respond point by point.

\medskip
\hrule
\medskip

\noindent \textbf{Response to the editor}

We are grateful for the editors’ constructive assessments
and have revised the manuscript accordingly.

\noindent \textbf{Editorial request:} ``Clarification of the significance of the tskit 1.0 release.''

We have revised the main text so that it now explicitly states that
the paper ``marks the release of tskit 1.0, a milestone that was chosen to formalise long-term
stability guarantees for its data formats and APIs.'' This is intended to make
clear that the significance of version 1.0 is the stabilisation of a mature
software platform for ARG analysis.

\noindent \textbf{Editorial request:} ``Contextualisation of the work with that
of others and your own prior work to clarify the advance made in the tskit
software.''

We have edited wording throughout to indicate some more of the context,
but due to space limitations, we have primarily addressed this with a new section
in the Supplementary Information. The new supplementary section then provides the
fuller developmental arc from the original tree sequence work, through
a series of publications developing core operations, to the present 1.0 release.

%% Previous version; I don't see how this answers the question, really?
% We have strengthened this in both the main text and the Supplementary
% Information. In the main text, we now explicitly state that, ``Building on the
% initial introduction of the succinct tree sequence data model and its
% formalisation as a general ARG representation, tskit 1.0 marks the maturity
% of the software library and data model for scalable ARG analysis (see
% Supplementary Information).''

\noindent \textbf{Editorial request:} ``Discussion comparing tskit to related
work with respect to functionality, efficiency, or scalability.''

To address this, we added the new supplementary section ``Tskit functionality
and benchmarks'' and Table S2. This material summarises the fundamental
algorithmic advances behind tskit's scalability, reviews published benchmarks, and
compares tskit with related libraries in terms of documented reusable
functionality.

\noindent \textbf{Editorial request:} ``Discussion of future directions.''

Given the word limit of the Correspondence format, we
addressed this concisely in the main text by adding the sentence: ``As tskit is
applied to a wider range of biological applications, future development is
likely to address additional complexities such as supporting multiple
chromosomes and structural variants.''

\medskip
\hrule
\medskip

\noindent \textbf{Response to Reviewer \#1}

\noindent \textbf{Reviewer \#1:} ``We strongly support the publication of this
announcement.''

We thank Reviewer 1 for this very supportive assessment, and especially for
taking the time to test the software directly. We are grateful for
the reviewer's recognition of tskit as enabling infrastructure for the field.

Reviewer 1 did not request any specific changes.

\medskip
\hrule
\medskip

\noindent \textbf{Response to Reviewer \#2}

We thank Reviewer 2 for the positive and thoughtful review, and for
highlighting the broad impact of tskit across simulation, inference, and
downstream population-genetic analysis. We are grateful for the reviewer’s
helpful suggestions regarding the significance of the 1.0 release and future
directions.

\noindent \textbf{Reviewer \#2:} ``Perhaps one comment is that the
correspondence could more clearly state the significance of the 1.0 release
(i.e., why this milestone is reached now).''

We have revised the manuscript so that the main text now explicitly
states that tskit 1.0 ``formalises long-term stability guarantees for its data
formats and APIs.'' The new supplementary section further explains that this
release marks the point at which a decade of development in representation,
algorithms, and reusable operations has been consolidated into a stable
software library.

\noindent \textbf{Reviewer \#2:} ``Also, it would be a valuable opportunity to
outline any future directions, and some perspective on how the tskit community
may evolve.''

Because of the strict word limit, we have addressed this concisely in
the main text by adding a sentence indicating likely future development
directions as tskit is applied to a broader range of biological applications.

\medskip
\hrule
\medskip

\noindent \textbf{Response to Reviewer \#3}

We thank Reviewer 3 for the careful and constructive review, and for
recognising the importance of software infrastructure for ARG-based methods. We
are grateful for the reviewer’s suggestions, which helped us clarify the
relationship to our earlier papers, better summarise the basis of tskit’s
scalability, and provide a more explicit comparison to related software.

\noindent \textbf{Reviewer \#3, comment 1:} ``The paper would benefit from
describing more clearly what is new in the present paper relative to the
original 2016 paper on an earlier version of tskit and a recent 2024 paper
describing the methodology. Is there anything fundamentally new beyond these
other publications? What does publishing this particular manuscript accomplish
that the 2024 paper did not? I do not think that needs to be explained in the
manuscript itself but might be better explained in a cover letter or response
to reviewers.''

We agree that this distinction needed to be made clearer. We have therefore
revised both the main text and the Supplementary Information to make the
developmental arc explicit. In brief, the 2016 paper introduced the succinct
tree sequence data model, the 2024 paper established it as a general
representation of ARGs, and the present Correspondence marks the maturity of
the software library itself as a stable platform for scalable ARG analysis. The
purpose of the present paper is therefore different from that of the earlier
methodological papers: it provides a concise and citable account of the mature
software platform, its stability guarantees, and the broad downstream ecosystem
that has formed around it.

\noindent \textbf{Reviewer \#3, comment 2:} ``Related to the prior point, it
would be helpful to have at least a summary of what makes this package
particularly effective for large-scale ARGs. Are there new innovations in
algorithms or data structures, for example? Or is it just better software
engineering?''

We have addressed this in the new supplementary section, which now
summarises the key algorithmic advances underlying
scalability, including the succinct tree sequence encoding, sequential tree
iteration, simplify, and later core operations for statistics, ancestry, and relatedness.
The section also points to published benchmarks demonstrating
major gains in speed and memory efficiency in several settings.

\noindent \textbf{Reviewer \#3, comment 3:} ``Related to the prior point, it
would be helpful to see more detail on comparison to related work, even if just
in the supplement, to back up the claims of superiority of tskit to
alternatives in the literature with respect to functionality, efficiency, or
scalability.''

We agree, and this was a major motivation for the new supplementary material.
We added the section ``Tskit functionality and benchmarks'' and Table S2, which
compare tskit with ARGneedle-lib, BTE/matUtils, and DendroPy across data model
features, file formats, metadata/provenance, tree operations, editing
operations, statistics, ancestry/relatedness, mutations/variants, and
visualisation. The accompanying text also explains the scoring methodology and
clarifies that the comparison is based on documented, reusable functionality
exposed through the Python APIs.

\medskip
\hrule
\medskip

We hope that the revised manuscript and Supplementary Information
satisfactorily address the editor's and reviewers' concerns. We are grateful
for the opportunity to revise the paper, and we believe the changes have
significantly improved the clarity and positioning of the manuscript.

The revised main text is now slightly over the word limit.
Given the positive overall assessment,
we were unsure what could most appropriately be cut
without weakening the manuscript. If the general direction and substance of these
revisions are viewed favourably, we would be very happy to shorten the text
further with editorial guidance on where compression would be most helpful.


\closing{Sincerely,}

\end{letter}
\end{document}
